\documentclass[../../../include/open-logic-section]{subfiles}

\begin{document}

\olfileid{his}{set}{pathology}
\olsection{اجیب اُساریاں}

پر پتا لگیا کہ \emph{حد} دی تعریف نال کجھ خاصیاں «اجیب» اُساریاں وی ممکن ہو جاندیاں نیں۔

انیسویں صدی دی چوتھی دہائی دے نیڑے بولزانو نے اک فنکشن لبھیا جیہڑا \emph{ہر تھاں مسلسل} سی پر \emph{کسے تھاں وی تفریق جوگ نہیں} سی۔ (بدقسمتی نال بولزانو نے ایہ گل کدے چھپوائی نہیں۔ ریاضی داناں نوں ایہ خیال پہلی واری 1872 وچ پتا لگیا، جدوں وائرشٹراس نے ایہو گل وکھرے طور اُتے لبھی۔)\footnote{ایہ تاریخ وکی پیڈیا دے \href{http://en.wikipedia.org/wiki/Weierstrass_function}{وائرشٹراس فنکشن} والے مضمون دیاں بہت تفصیلی حاشیاں وچ دَسی گئی اے۔} گھٹ توں گھٹ اینا تاں آکھیا جا سکدا اے کہ ایہ بڑی حیرانی والی گل سی۔ $|x|$ ورگے فنکشن لبھنے سوکھے نیں جیہڑے ہر تھاں مسلسل نیں پر کسے خاص نقطے اُتے تفریق جوگ نہیں۔ پر جیہڑا فنکشن ہر تھاں مسلسل ہووے تے \emph{کسے تھاں وی} تفریق جوگ نہ ہووے، اوہ بالکل وکھری گل اے۔ ذرا سوچو ایہو جیہا فنکشن کِویں کھچو گے۔ تسلسل لئی لازم اے کہ کاغذ توں قلم چُکے بغیر سارا گراف کھچیا جا سکے۔ پر کسے تھاں وی تفریق جوگ نہ ہون لئی قلم دی سمت لگاتار تے اچانک بدلنی پئے گی۔

ہور «اجیب» مثالاں وی سامنے آئیاں۔ 5 جنوری 1874 نوں کانتور نے ڈیڈیکائنڈ نوں خط وچ ایہ سوال پُچھیا:
\begin{quote}
کی کسے سطح (مثلاً چُفیرے دی حد سمیت اک مربع) نوں کسے لکیر (مثلاً دوہاں سِریاں سمیت اک سدھی لکیر) نال اک بہ اک جوڑیا جا سکدا اے، ایویں کہ سطح دے ہر نقطے نال لکیر دا اک نقطہ تے، اُلٹے پاسے، لکیر دے ہر نقطے نال سطح دا اک نقطہ جُڑے؟

ایس ویلے وی مینوں لگدا اے کہ ایس سوال دا جواب بڑا اوکھا اے۔ بھاویں ایتھے وی \emph{نہ} آکھن نوں اینا چِت کردا اے کہ ثبوت لگ بھگ وادھو ای لگے۔ [حوالہ \citealt{Gouvea2011} وچ اے۔]
\end{quote}
پر 1877 وچ کانتور نے ثابت کیتا کہ اوہ غلط سی۔ لکیر تے مربع وچ نقطیاں دی گنتی عین اِکّو اے۔ اوہنے 29 جون 1877 نوں ڈیڈیکائنڈ نوں لکھیا، «\emph{je le vois, mais je ne le crois pas}»، یعنی «میں ایہ ویکھدا آں، پر مینوں یقین نہیں آؤندا۔» ریاضی دی «رائج تاریخ» وچ ایس گل توں اکثر ایہ مطلب لیا جاندا اے کہ اوس ویلے دے ریاضی داناں نوں ایہ نویں نتیجے \emph{لفظی معنیاں وچ ناقابلِ یقین} لگے سن۔ (خطاں دی گل بات \citet{Gouvea2011} وچ ملدی اے، تے اسیں \olref[his][set][mythology]{sec} وچ اوس ول مُڑاں گے۔ کانتور دے ثبوت دا خاکہ \olref[his][set][cantorplane]{sec} وچ اے۔)

کانتور دے نتیجے توں تحریک لے کے پیانو ایہ سوچن لگ پیا کہ کی لکیر نوں کسے مستوی اُتے \emph{ہمواری نال} بھیجیا جا سکدا اے۔ ایہ اک \emph{ساری تھاں بھرن والی منحنی} ہووے گی۔ \citeyear{Peano1890} وچ پیانو نے ایہو جیہی منحنی بنائی۔ ایہ سچ مُچ عام سہج سمجھ دے اُلٹ اے۔ اقلیدس نے لکیر نوں «چوڑائی توں بناں لمبائی» آکھیا سی (پہلی کتاب، دوجی تعریف)، پر پیانو نے وکھا دِتا کہ لکیر نوں ڈھکدی طرح وَل دے کے اوہدی لمبائی توں چوڑائی بنائی جا سکدی اے۔ \citeyear{Hilbert1891} وچ ہلبرٹ نے سمجھ وچ ذرا سوکھی آؤن والی ساری تھاں بھرن والی منحنی دَسی تے اوہنوں وکھاؤن لئی کجھ تصویراں وی دتیاں۔ منحنی اک ترتیب نال اُسردی اے۔ اُساری دے پہلے چھے پڑا ایہ نیں:
\begin{center}
\begin{tikzpicture}
\begin{scope}[xshift=20pt, yshift=20pt]
	\draw [oldiagcolorC, l-system={Hilbert curve, step=40pt, angle=90, axiom=L, order=1}] lindenmayer system;
\end{scope}
\begin{scope}[xshift=110pt, yshift=10pt]
	\draw [oldiagcolorC, l-system={Hilbert curve, step=20pt, angle=90, axiom=L, order=2}] lindenmayer system;
\end{scope}
\begin{scope}[xshift=205pt, yshift=5pt]
	\draw [oldiagcolorC, l-system={Hilbert curve, step=10pt, angle=90, axiom=L, order=3}] lindenmayer system;
\end{scope}
\begin{scope}[xshift=2.5pt, yshift=-97.5pt]
	\draw [oldiagcolorC, l-system={Hilbert curve, step=5pt, angle=90, axiom=L, order=4}] lindenmayer system;
\end{scope}
\begin{scope}[xshift=101.25pt, yshift=-98.75pt]
	\draw [oldiagcolorC, l-system={Hilbert curve, step=2.5pt, angle=90, axiom=L, order=5}] lindenmayer system;
\end{scope}
\begin{scope}[xshift=200.625pt, yshift=-99.375pt]
	\draw [oldiagcolorC, l-system={Hilbert curve, step=1.25pt, angle=90, axiom=L, order=6}] lindenmayer system;
\end{scope}
\draw[gray] (0pt,0pt) rectangle (80pt, 80pt);
\draw[gray] (100pt,0pt) rectangle (180pt, 80pt);
\draw[gray] (200pt,0pt) rectangle (280pt, 80pt);
\draw[gray] (0pt,-100pt) rectangle (80pt, -20pt);
\draw[gray] (100pt,-100pt) rectangle (180pt, -20pt);
\draw[gray] (200pt,-100pt) rectangle (280pt, -20pt);
\end{tikzpicture}  
\end{center}
حد وچ، جیہدی ہُن دقیق تعریف بن چکی سی، سارا مربع !!{colorC} رنگ نال پُورا بھر جاندا اے۔ نال ایہ وی دَس دیئے کہ ہلبرٹ دی منحنی ہر تھاں مسلسل اے پر کسے تھاں وی تفریق جوگ نہیں۔ عام سمجھ نال ایہ ایویں اے کہ لامتناہی حد وچ فنکشن ہر پل اچانک پاسا بدلدا اے۔ (ہلبرٹ دی اُساری دا ہور تفصیلی خاکہ اسیں \olref[his][set][hilbertcurve]{sec} وچ دیاں گے۔)

چنگا ہویا یا ماڑا، ایہناں «اجیب» ہندسی اُساریاں نوں ہندسے دی سہج سمجھ اُتے بھروسا کرن دے خلاف اک وجہ سمجھیا گیا۔ ایہ ریاضی کرن دے اک نویں طریقے دا لگ بھگ \emph{پرچار} بن گئیاں، جیہڑا اگے چل کے سیٹ نظریے تک پہنچیا۔ ایس موضوع دیاں بعد وچ گھڑیاں گئیاں متھاں وچ وار وار دُہرایا گیا کہ ایہ نتیجے پوری طرح دقیق تے بالکل حیران کر دین والے سن۔ ایس طرح ایہناں دا دوہرا کم سی۔ ہندسی سہج سمجھ اُتے بھروسا کرن توں خبردار کرنا تے وکھاؤنا کہ نویں ڈھنگ نال سوچن توں کِنّے نتیجے نکل سکدے نیں۔

\end{document}
